\documentclass[11pt]{article}
% Draft of the third bridges article. Built with XeLaTeX:
% cd docs/dualities && tectonic article.tex
\usepackage{polyglossia}\setmainlanguage{english}\setotherlanguage{french}
\newcommand{\thmname}{Theorem}\newcommand{\corname}{Corollary}\newcommand{\propname}{Proposition}
\newcommand{\lemname}{Lemma}\newcommand{\bridgename}{Bridge}\newcommand{\defname}{Definition}
\newcommand{\exname}{Example}\newcommand{\exoname}{Exercise}\newcommand{\remname}{Remark}
% Preamble of docs/dualities/article.tex: the bridges preamble, with fonts loaded by file name
% so that tectonic builds it.
% Built with XeLaTeX: cd docs/bridges && xelatex article.tex && xelatex article.tex
\usepackage{fontspec}
\setmainfont{XCharter}[Extension=.otf, UprightFont=*-Roman, BoldFont=*-Bold, ItalicFont=*-Italic, BoldItalicFont=*-BoldItalic]
\setmonofont{DejaVuSansMono}[Extension=.ttf, BoldFont=*-Bold, ItalicFont=*-Oblique, BoldItalicFont=*-BoldOblique, Scale=0.8]
\usepackage{microtype}
\usepackage[a4paper,margin=27mm]{geometry}
\usepackage{setspace}\setstretch{1.1}
\usepackage{amsmath,amssymb,amsthm}
\usepackage{mathtools}
\usepackage[most]{tcolorbox}
\usepackage{xcolor}
\usepackage{tikz}
\usetikzlibrary{arrows.meta,positioning,calc,cd}
\usepackage{booktabs,array,tabularx}
\usepackage{fancyvrb}
\usepackage[hidelinks]{hyperref}
\usepackage{xurl}
\usepackage{enumitem}
\setlist{nosep,leftmargin=*}
\setlength{\emergencystretch}{2em}
\setlength{\parskip}{0.35em}
\numberwithin{equation}{section}

\newcommand{\Set}{\mathbf{Set}}
\newcommand{\PSh}{\mathrm{PSh}}
\newcommand{\Sh}{\mathrm{Sh}}
\newcommand{\Kar}{\operatorname{Kar}}
\newcommand{\Hom}{\operatorname{Hom}}
\newcommand{\End}{\operatorname{End}}
\newcommand{\Aut}{\operatorname{Aut}}
\newcommand{\id}{\mathrm{id}}
\newcommand{\B}{\mathbf{B}}
\newcommand{\y}{\mathbf{y}}
\newcommand{\op}{\mathrm{op}}
\newcommand{\res}[2]{#1|_{#2}}
\newcommand{\q}[1]{«\,#1\,»}
\newcommand{\lean}[1]{{\small\nolinkurl{#1}}}
% let long declaration names break at underscores and dots
\makeatletter\g@addto@macro\UrlBreaks{\do\_\do\.}\makeatother
\newcommand{\lp}{\href{https://github.com/LeanPool/lean-pool/tree/main/LeanPool/ToposTheoreticBridges}{\nolinkurl{ToposTheoreticBridges}}}

% Lean code blocks.
\DefineVerbatimEnvironment{leancode}{Verbatim}{fontsize=\small,frame=leftline,framerule=0.6pt,rulecolor=\color{blue!40},xleftmargin=6pt}

\theoremstyle{plain}
\newtheorem{theorem}{\thmname}[section]
\newtheorem{corollary}[theorem]{\corname}
\newtheorem{proposition}[theorem]{\propname}
\newtheorem{lemma}[theorem]{\lemname}
\newtheorem{bridge}{\bridgename}
\renewcommand{\thebridge}{\Alph{bridge}}
\theoremstyle{definition}
\newtheorem{definition}[theorem]{\defname}
\newtheorem{example}[theorem]{\exname}
\newtheorem{exercise}{\exoname}[section]
\theoremstyle{remark}
\newtheorem*{remark}{\remname}
\tcolorboxenvironment{theorem}{enhanced, breakable, colback=green!2, colframe=green!35!black, boxrule=0.6pt, arc=0pt, left=7pt, right=7pt, top=4pt, bottom=4pt, before skip=12pt, after skip=8pt}
\tcolorboxenvironment{bridge}{enhanced, breakable, colback=blue!3, colframe=blue!35!black, boxrule=0.5pt, arc=0pt, left=7pt, right=7pt, top=4pt, bottom=4pt, before skip=12pt, after skip=8pt}
\tcolorboxenvironment{definition}{enhanced, breakable, colback=black!3, colframe=black!25, boxrule=0.4pt, arc=0pt, left=7pt, right=7pt, top=4pt, bottom=4pt, before skip=10pt, after skip=6pt}
\newcommand{\entry}[1]{\par\smallskip\noindent\textbf{#1}\enspace}

% Pipeline trace of a generated theorem (labels set per language in article*.tex).
\newenvironment{trace}{\par\begin{tcolorbox}[enhanced,breakable,colback=yellow!4,colframe=yellow!45!black,boxrule=0.4pt,arc=0pt,fontupper=\small,left=5pt,right=5pt,top=3pt,bottom=3pt,before skip=8pt,after skip=8pt]\textbf{\tracename}\par\begin{description}[leftmargin=2.6cm,labelwidth=2.4cm,labelsep=0.2cm,font=\normalfont\itshape]}{\end{description}\end{tcolorbox}}

\newcommand{\Id}{\mathrm{Id}}

\title{Stone-type dualities from one bridge\\[0.3em] \large Running Caramello's duality machine mechanically, and checking its outputs against Mathlib}
\author{Michel Hua\\ \small Commutator, Paris\\ \small\texttt{michel@commutator.io} · ORCID 0009-0002-8304-0683}
\date{Draft outline, 28 November 2026}

\begin{document}
\maketitle

\begin{abstract}
Caramello showed that the classical Stone-type dualities are instances of one topos-theoretic
phenomenon \cite{caramello-stone}: the toposes of sheaves on a preorder with a coverage and on
its locale of ideals are equivalent, and a duality is what one reads when a structure is
recovered from that topos through an invariant. We run this machine mechanically. Its core, the
Comparison Lemma for preorders \cite[Thm.~10.13(i)]{caramello-stone}, is proved once in Lean~4,
without the axiom of choice. Each duality is then the choice of three inputs (a preorder, a
coverage, a dense part) and two readings, one on each side, and its output is compared with the
statement already in Mathlib. Birkhoff's representation of finite distributive lattices comes out
as Mathlib's \lean{OrderIso.lowerSetSupIrred}, definitionally, and Lindenbaum--Tarski duality as
the isomorphism of a complete atomic Boolean algebra with the power set of its atoms. We list the
dualities of Mathlib that the machine reaches, what each one needs beyond the machine (a choice of
points, and for Stone duality a form of choice), and those it does not reach: Gelfand,
Pontryagin, Morita theory of rings, Tannaka duality. The text is a draft; its sections are
outlined below, with the parts already checked in Lean marked as such.
\end{abstract}

\section{The question}\label{sec:intro}
The companion article \cite{bridges-generator} runs Caramello's generator on presheaf toposes,
and recovers two classical theorems at once in its section on equivalences: Birkhoff's theorem
for finite distributive lattices and the correspondence between finite $T_0$ spaces and finite
posets. Here we ask how far the same mechanism goes for the dualities that are already proved in
Mathlib, when it is run without any idea added by hand: which inputs it needs, which readings,
and what has to be supplied that is not a mechanical step. The interest is not in the dualities,
which are classical, but in the boundary between what the machine produces and what it does not.

\section{The machine}\label{sec:machine}
\paragraph{Inputs.} Following \cite[\S3, \S9]{caramello-stone}: a preorder $C$; a coverage $J$ on
$C$ given by covering sieves; and a $J$-dense part $D\subseteq C$, meaning that every $c$ is covered
by the sieve of elements lying below some $d\in D$ with $d\le c$. For a duality with spaces rather
than locales, a fourth input: a set of points of the topos.

\paragraph{The bridge.} The $J$-ideals of $C$ (lower sets that contain every element covered by a
sieve inside them) form the locale $\Id_J(C)$, and $\Sh(C,J)\simeq\Sh(\Id_J(C))$
\cite[Thm.~3.1]{caramello-stone}. Its order-theoretic core is the Comparison Lemma for preorders:

\begin{leancode}
def PreorderCoverage.densityOrderIso (hD : J.Dense D) :
    J.Ideal ≃o (J.restrict D).Ideal
\end{leancode}
Intersecting with $D$ is an order isomorphism from the $J$-ideals of $C$ to the ideals of the
coverage that $J$ induces on $D$, with inverse the $J$-closure. Lean: \lean{Bridges/Dualities/Density.lean};
\verb|#print axioms| gives \lean{propext} and \lean{Quot.sound}, no choice, in line with
Caramello's remark that the dualities with locales are constructive.

\paragraph{Readings.} A duality is obtained by reading the two sides of the isomorphism: on the
left, the $J$-ideals of $C$ in terms of $C$ (for the coverages used below, they are the principal
ideals, so $C$ itself); on the right, the induced coverage on $D$ (often trivial, so its ideals
are the lower sets of $D$, or all its subsets).

\section{Outputs checked in Lean}\label{sec:outputs}

\begin{bridge}[Lindenbaum--Tarski]
A complete atomic Boolean algebra $B$ is isomorphic to the power set of its atoms.
\end{bridge}
\emph{Inputs:} $C=B$; the canonical coverage (a lower set below $a$ covers $a$ when its join is
$a$; stability uses the frame law $b\wedge\bigvee S=\bigvee_{s}b\wedge s$); $D$ the atoms, dense
because $B$ is atomistic. \emph{Readings:} the canonical ideals are the principal ideals; on the
atoms the induced coverage is trivial (a covering sieve of an atom contains it), so every set of
atoms is an ideal. Lean: \lean{completeAtomicBooleanAlgebraOrderIsoSetAtoms}, sending $a$ to the
atoms below it.

\begin{bridge}[Birkhoff, finite]
A finite distributive lattice $L$ is isomorphic to the lattice of lower sets of its
sup-irreducible elements.
\end{bridge}
\emph{Inputs:} $C=L$; the coherent coverage (finite joins; stability uses distributivity); $D$ the
sup-irreducible elements, dense by \lean{exists_supIrred_decomposition}. \emph{Readings:} on a
finite lattice the coherent ideals are principal; on the sup-irreducibles the induced coverage is
trivial (\lean{SupIrred.finset_sup_eq}), so its ideals are the lower sets. Lean:
\lean{birkhoffOrderIso}, and \lean{birkhoffOrderIso_apply} proves it equal to Mathlib's
\lean{OrderIso.lowerSetSupIrred}.

\section{The registry of targets}\label{sec:registry}
\begin{center}\small
\begin{tabularx}{\linewidth}{>{\raggedright\arraybackslash}p{3.6cm}>{\raggedright\arraybackslash}X>{\raggedright\arraybackslash}p{2.6cm}}
\toprule
duality (Mathlib) & inputs of the machine & status\\
\midrule
Lindenbaum--Tarski & complete atomic Boolean algebra, canonical coverage, atoms & checked in Lean\\
Birkhoff, finite (\lean{Order/Birkhoff.lean}) & finite distributive lattice, coherent coverage, sup-irreducibles & checked, equal to Mathlib's\\
preorders and Alexandrov spaces (\lean{AlexandrovDiscrete.lean}) & a preorder, trivial coverage; points: the principal filters & to do\\
spatial frames and sober spaces (\lean{FrameAdjunction.lean}) & a frame, canonical coverage; points: completely prime filters \cite[\S4.3]{caramello-stone} & to do\\
Stone, distributive lattices and Boolean algebras & coherent coverage; points: prime filters, whose existence needs the prime ideal theorem \cite[\S4.2]{caramello-stone} & to do; the choice step is not mechanical\\
Priestley, Esakia (\lean{Topology/Order/Priestley.lean}) & as Stone, with the order on prime filters & to do; Mathlib has the spaces, not the duality\\
Galois theory, coverings & atomic toposes and their points \cite{caramello-galois} & outside the preorder case; later\\
Gelfand, Pontryagin, Morita of rings, Wedderburn--Artin, Tannaka & no preorder site; not toposes of this kind & out of reach (\S\ref{sec:limits})\\
\bottomrule
\end{tabularx}
\end{center}

\section{What is mechanical, and what is not}\label{sec:mechanical}
To be written from the runs: the Comparison Lemma and the two readings are uniform, and each
reading is a short lemma about the chosen coverage; the choices of coverage and of dense part are
where the duality is decided, and they come from the literature, not from the machine; the passage
to spaces needs enough points, which for Stone duality is the prime ideal theorem.

\section{Out of reach}\label{sec:limits}
To be written: why Gelfand and Pontryagin duality, Morita theory of rings and Tannaka duality are
not outputs of this machine (no site whose topos carries them; categories of modules are not
toposes), and what a bridge in the broader sense would need there.

\begin{thebibliography}{9}
\bibitem{caramello-stone} O.~Caramello, \emph{A topos-theoretic approach to Stone-type dualities},
arXiv:1103.3493, 2011.
\bibitem{caramello-galois} O.~Caramello, \emph{Topological Galois theory}, arXiv:1301.0300, 2013.
\bibitem{bridges-generator} M.~Hua, \emph{Topos-theoretic bridges as a theorem generator},
draft, docs/bridges, 2026.
\end{thebibliography}
\end{document}
